\section{A parameter-free MP command for the PA196 calculus}
\label{sec:residual-pa}

This construction transfers the opposite-comparison-residual idea from the
repository's ZFC \texttt{k215-residual-parameter} experiment to PA196's
different three-command language. It is a source-level construction. No
claim of a smaller physical machine or a newly kernel-checked PA theorem is
made here.

Write $[a]W=P(a,W)$ for the exact unshifted Cantor stack constructor,
$I(a,b)=4P(a,b)$, and $U(i,a)=4P(a,i)+2$. Code $0$ denotes truth and code
$1$ is the target contradiction. A store is $(t,W,d,h)$. Retain Shift,
which halts if $t+h=1$ and otherwise maps
\[
 S(t,[a]W,d,h)=(h,W,d,a).
\]
Replace the parameterized MP command by the nullary command
\[
 M(t,W,d,h)=
 \begin{cases}(d,[0]W,d,0)& I(t,d)=h,\\(0,W,d,0)&I(t,d)\ne h.
 \end{cases}
\]
For $A(c,x,y)$ put $D=d+x$, $H=h+y$. Emit exactly the old gate
$g(H,D,c)$ and exactly the old eight heads $r_1,\ldots,r_8$, in their
old top-to-bottom order; in particular $r_8=U(d,h)$ and $r_7$ is the
restricted implicational head with fourth parameter $d\le D$.
Change only the numeric residual:
\[
 A(c,x,y)(t,W,d,h)=
 \bigl(g,[r_1]\cdots[r_8][t]W,\max(c-D-1,0),0\bigr).
\]
Every emitted head and gate is different from $1$, and $r_8\ge2$.
These are literal constructor-tag facts, requiring no assumption that
$t,h$, or the old stack are sound.

\subsection{Exact command macros}

Commands in a displayed word execute from left to right. Let $A_0=A(0,0,0)$
and define
\[
 P=A_0S^8M,\qquad R=A_0S^{10},\qquad
 Z=(PR)^2,\qquad
 T_v=P\,A(v+1,0,v+1)S^{10}.
\]
The following equations hold for arbitrary natural stores, with no early
halt:
\begin{align*}
 P(t,W,d,h)&=(0,[t]W,0,0),\\
 R(t,[a]W,d,h)&=(t,W,0,a),\\
 Z(t,W,d,h)&=(0,W,0,0),\\
 T_v(t,W,d,h)&=(0,W,v,t).
\end{align*}
For $P$, $A_0$ sets the residual to zero. After eight Shifts the cache is
$U(d,h)$, of residue $2$ modulo $4$; hence $M$ fails regardless of the
current value and returns zero without touching the saved old top.
All tested adjacent fresh values differ from $1$, so their sum cannot be
$1$. For $R$, the tenth Shift tests $U(d,h)+t\ge2$, then exposes $t$ and
pops one old stack item. Applying $PR$ twice proves the formula for $Z$.
After $P$ in $T_v$, both theorem fields and $d$ are zero. The second AX
has $D=0$, $H=c=v+1$, so its gate is zero and its residual is $v$. Its
generalization is $U(0,0)=2$. Ten Shifts expose the stored original $t$
in the cache and restore the original stack.

These equations use only $\operatorname{decode}(P(a,W))=(a,W)$.
They therefore remain valid when $a=W=0$ and $P(0,0)=0$; no injective
distinction between an empty stack and a pushed zero is presumed.

\subsection{A proof compiler}

A compiler block for a derivable formula $B$ is required, on any numerical
store, either to halt or to return $(B,W,d',0)$, preserving the old stack
exactly. The parameter $d'$ may be arbitrary. The following construction is
structural induction on a derivation in PA196's exact numeric calculus.

Truth is compiled by $Z$. To select the $j$th AX head, $1\le j\le8$,
start AX with current zero and follow it by
\[
 S^{j+1}R^{8-j}.
\]
The initial Shifts place $r_j$ in the current field. Subsequent $R$s
preserve that field and discard only the remaining fresh prefix, ending
with the saved zero in the cache. For the restricted head with data
$(D,H,c,u)$, $u\le D$, the complete block is
\[
 Z\,T_u\,A(c,D-u,H)\,S^8R.
\]
The other six heads use $u=0$ and the appropriate selection index.
For a permitted gate, use $Z A(c,D,H) R^9$. The nine pops remove the
eight new heads and saved zero while retaining the gate as current.

If $p_B$ compiles $B$, a block for $U(i,B)$ is
\[
 p_B\,T_i\,A_0S^9.
\]
After $T_i$ the store is $(0,W,i,B)$; the last eight-head item is exactly
$U(i,B)$, and nine Shifts expose it with cache zero.

If $p_I$ compiles $I(A,B)$ and $p_A$ compiles $A$, modus ponens is compiled
by
\[
 p_I\,P\,p_A\,T_B\,S\,M\,R.
\]
Unless a premise block has already halted, the successive key stores are
\[
 (I(A,B),W,d,0),\quad(0,[I(A,B)]W,0,0),\quad
 (A,[I(A,B)]W,d_1,0),\quad(0,[I(A,B)]W,B,A).
\]
The Shift halts when $A=1$, which is acceptable. Otherwise it gives
$(A,W,B,I(A,B))$, so $M$ succeeds and gives $(B,[0]W,B,0)$.
The final $R$ removes precisely that zero, yielding $(B,W,0,0)$.
This also covers $I(0,0)=0$ literally. For a derivation of code $1$,
append one Shift to its compiler block: $(1,W,d',0)$ passes the exact
halt test.

Soundness uses the usual invariant: the current value, cached value, and
every numeric stack entry are derivable. Shift can halt only when one of
the two theorem values is $1$ and the other is $0$. Successful MP is
literal modus ponens, failed MP returns truth, and the altered AX emits
the same formulas as before. Its numeric residual has no logical invariant.
Thus the revised command calculus searches for exactly the same refutations.

\subsection{Changed arities and fair enumeration}

Let $e(\epsilon)=0$ and encode a nonempty continuation by
\begin{align*}
 e(Sw)&=1+P(0,e(w)),\\
 e(Mw)&=1+P(1,e(w)),\\
 e(A(c,x,y)w)&=1+P(c+2,P(y,P(x,e(w)))).
\end{align*}
These are the old tags; only MP no longer consumes an additional operand
field. The untouched scheduler subtracts one from positive continuations
and decodes the outer pair. Deleting its old 14-site reader implements
exactly this new convention. Each positive continuation decreases under
decoding, so a finite word is consumed completely. This is a claim about
the rewritten source, not an assertion that old certificate integers
have their old meanings.

The distinction between a continuation and a fresh candidate is essential.
For nonempty $w$, put $f(w)=e(w)-1$. If the continuation register is
$L=e(w)>0$, the scheduler decodes $L-1=f(w)$ and leaves its fresh
candidate counter $N$ unchanged. If $L=0$ and $N=f(w)$, it decodes $N$
without subtracting one and increments that counter. Both cases therefore
decode the same first command and leave $L=e(w')$ after its parameter
fields, where $w'$ is the tail. This is a structural induction on $w$,
using the exact Cantor inverse identities. The original source supplies
precisely these two interfaces as \texttt{Fetch.positive} and
\texttt{Fetch.fresh} in \texttt{PA196/Source/Fetch.lean}; its
\texttt{FetchDispatch.seed} in \texttt{Source/FetchDispatch.lean} is
literally $N$ when $L=0$, and $L-1$ otherwise. The source splice changes
neither of those routines. Parameter-free MP consumes no field after its
tag, whereas AX still consumes $y$ and then $x$.

The compiler contract is uniform in the incoming numerical store, hence
cofinality is obtained by prefixing a fixed compiled refutation with
$A(Y,0,0)Z$. If $z$ is the code of the fixed suffix, the fresh-candidate
integer is $P(Y+2,P(0,P(0,z)))$. It is strictly increasing and unbounded
in $Y$, so an increasing candidate enumeration beginning at $2$ includes
such a halting word. Conversely, every halt from a theorem-safe seed is
a derivation of $1$. The needed coding and block arguments use finite
induction and primitive-recursive proof transformations; a PA-internal
binding to a literal TM still requires the separate machine formalization.

\subsection{Exact source splice and its present cost}

In the pinned PA196 source, replace the cleanup subtree
$({\rm CD}6,{\rm CC}2)$ at old sites $3687$--$3688$ by
$({\rm CD}6,({\rm CI}4,{\rm CC}2))$ within the second comparison
$[3680,3692)$. The paired comparison loop consumes one excess scratch
unit on inequality. Its remaining scratch value is
$\max(c-D-1,0)$, transferred to initially zero register $4$; registers
$5,6$ end zero. Then remove the closed MP reader $[3693,3707)$.
Every unchanged instruction and both of its ordinal successors are
checked under the explicit splice map by \texttt{residual\_source.py}.
Two independent CFG resolvers agree. The isolated comparison was also
replayed for all $0\le c,D<64$; this finite replay supplements the loop
argument and does not replace it.

Here is the complete arbitrary-natural comparison argument. At its entry
the registers $(T,H,X,d)$ are $(t,D,c,0)$. Number the thirteen new
instructions $0,\ldots,12$ and call the exit $13$. Their resolved effects
are
\[
\begin{array}{c|l@{\quad}c|l}
0&{\rm dec}\ H:(4,1)&1&{\rm dec}\ X:(3,2)\\
2&{\rm jump}\ 0&3&{\rm jump}\ 5\\
4&{\rm dec}\ X:(11,5)&5&{\rm dec}\ H:(7,6)\\
6&{\rm jump}\ 5&7&{\rm dec}\ X:(10,8)\\
8&{\rm inc}\ d;9&9&{\rm jump}\ 7\\
10&{\rm jump}\ 13&11&{\rm dec}\ T:(13,12)\\
12&{\rm jump}\ 11&&
\end{array}
\]
where the two targets of a decrement are zero and positive, respectively.
After $k\le\min(D,c)$ complete trips through $0,1,2$, the invariant is
$(T,H,X,d)=(t,D-k,c-k,0)$ and the control is $0$. Each trip takes three
steps. The terminal cases are:
\begin{itemize}
\item If $D=c$, rows $0,4$ enter $11$ with $H=X=0$.
The $11,12$ loop takes $2t+1$ steps and clears $T$; $d$ remains zero.
\item If $D<c$, rows $0,4$ enter $5$ with
$(H,X,d)=(0,c-D-1,0)$, leaving $T$ untouched.
\item If $D>c$, rows $0,1,3$ enter $5$ with
$(H,X,d)=(D-c-1,0,0)$, again leaving $T$ untouched.
\end{itemize}
Rows $5,6$ clear $H$ in $2H+1$ steps. In rows $7,8,9$, if the incoming
$X$ is $r$, after $j\le r$ trips the invariant is
$(X,d)=(r-j,j)$; each trip takes three steps. The zero test and exit
jump take two more steps. Thus the exact comparison duration is
\[
 \begin{cases}
 3D+2t+3,&D=c,\\
 3c+2,&D<c,\\
 2D+c+4,&D>c.
 \end{cases}
\]
The final store is always
$(T,H,X,d)=({\rm if}\ D=c\ {\rm then}\ 0\ {\rm else}\ t,0,0,
\max(c-D-1,0))$. Every other register is unchanged. This proves
termination and the residual contract by elementary natural-number
induction, including both zero boundary cases. The preceding unchanged
guard copies $c$ twice, performs the first $H=c$ test, then moves old
$D$ to $H$, leaving the residual destination $d=0$ exactly as required.
Consequently the final gate tests are still the two original equalities;
only the disposable excess $X$ is now retained as a numeric parameter.

The residual transfer has \emph{exactly the same register-machine execution
time} as the old comparison, for every $t,D,c$. In the old cleanup, one
positive $X$ cycle consists of decrementing $X$, jumping back to the
$H$ test, and testing the already zero $H$ before returning to the $X$
test: three steps. In the new cleanup, the corresponding cycle decrements
$X$, increments $d$, and jumps directly back to the $X$ test: also three
steps. The initial $H$ drain and both exit tests are unchanged. The shared
prefix, the equality branch, and the $D>c$ branch have unchanged step
counts as well. Applying the two cycle invariants to arbitrary $r$ proves
that both comparisons obey the three duration formulas above, not just
that a finite selection of runs agrees. The constructor additionally
checks both old and new outputs and equal durations in all 4,096 probes.
Thus the numeric information was available in an otherwise wasted loop
test. This does not imply equal physical TM time: the new source address
layout still has to be compiled and verified.

The result has $3,799$ source instructions and $4,941$ guarded cases.
It has $164$ distinct internal subtrees versus $159$ in PA196 because
the changed comparison loses existing sharing. The source SHA-256 is
\texttt{9aa46652ee8f11d7919700f860578d3c4c6b0ad8ec43524af79ec7d4703dff4a}.
A conservative startup-preserving allocator with no new labels reports
an empty domain. This is a limitation of that allocation model, not an
impossibility theorem. No physical state reduction is established.
